% p3-appC-notation

\section{P3 Appendix C: Notation Summary}

Appendix C: Notation Summary

  Figure (fig-p3-appC-notation). Notation-summary triptych --- Glyph table / Operators / Constants.

  - $\varphi$: golden ratio (1+5)/2 approx 1.61803.
  - psi: reciprocal golden ratio $\varphi$-1=$\varphi$-1approx0.61803.
  - (star): Trinity anchor identity $\varphi$2+$\varphi$-2=3.

  - Z[$\varphi$]: Lucas ring {a+b$\varphi$:a,b$\in$Z}.
  - T(n,k,p,m,q): Trinity monomial n$\cdot$3k$\cdot\varphi$p$\cdot$pim$\cdot$ eq.
  - lambda: Trinity label (n,k,p,m,q).
  - ell(lambda): bit-complexity of a Trinity label.
  - T: set of all Trinity monomials.
  - SL: complexity shell at level L.
  - PiL(x): projection of x to a complexity-L Trinity approximant.
  - RL(x): residual x-PiL(x).
  - sigma, K, Lw: Pellis shift, Koopman operator, and transfer operator.
  - MDL(T\textasciicircum{}*,c): MDL score of the unification representation.
  - muT(x): Trinity approximation exponent.
  - mu\textasciicircum{}*=1: critical exponent of the symbolic phase transition.
  - M$\in$R42: catalogue MDL vector.
  - alpha$\varphi$( alg): algebraic AlphaPhi, $\varphi$-3/2approx0.118.
  - alpha$\varphi$(log): logarithmic AlphaPhi, 2ln$\varphi$/piapprox0.30635.
  - Companion Paper 1: Pellis hierarchical $\varphi$-expansions.
  - Companion Paper 2: Trinity monomial representations.